\documentclass[a4paper,11pt]{article}
\usepackage[utf8]{inputenc}
\usepackage[french]{babel}
\usepackage[T1]{fontenc}
\usepackage{amsmath,amssymb,amsthm}
\usepackage{geometry}
\geometry{margin=2cm}
\usepackage{enumitem}
\usepackage{hyperref}
\usepackage{booktabs}
\usepackage{array}
\usepackage{xcolor}
\usepackage{listings}
\usepackage{titlesec}
\usepackage{fancyhdr}
\usepackage{lastpage}
\usepackage{graphicx}
\usepackage{etoolbox}
\AtBeginDocument{\shorthandoff{;:!?}}

\titleformat{\section}{\large\bfseries\color{blue!60!black}}{}{0pt}{}[\vspace{-0.5ex}\rule{\textwidth}{0.4pt}]
\titleformat{\subsection}{\bfseries\color{blue!40!black}}{}{0pt}{}
\titleformat{\subsubsection}{\itshape\color{blue!30!black}}{}{0pt}{}

\lstdefinestyle{pystyle}{
  frame=single,
  numbers=left,
  numbersep=5pt,
  breaklines=true,
  basicstyle=\small\ttfamily,
  backgroundcolor=\color{gray!5},
  rulecolor=\color{gray!40},
  columns=fullflexible,
  keepspaces=true,
  showstringspaces=false,
  language=Python
}
\lstset{style=pystyle}

\pagestyle{fancy}
\fancyhf{}
\fancyhead[L]{\small STA211 -- Stacking (Python)}
\fancyhead[R]{\small Fiche r\'ecapitulative}
\fancyfoot[C]{\thepage/\pageref{LastPage}}
\renewcommand{\headrulewidth}{0.4pt}

\theoremstyle{definition}
\newtheorem*{definition}{D\'efinition}
\theoremstyle{remark}
\newtheorem*{remark}{Remarque}
\newtheorem*{important}{\`A retenir}

\title{\textbf{Fiche r\'ecapitulative -- STA211}\\[4pt]
\large Stacking : agr\'egation par m\'eta-apprentissage \\[2pt]
\small Impl\'ementation Python avec scikit-learn}
\author{Vincent Audigier, Nd\`eye Niang (CNAM)}
\date{15 juin 2026}

\begin{document}
\maketitle
\thispagestyle{empty}

\begin{abstract}
Le \textbf{stacking} (ou \textit{stacked generalization}) est une
m\'ethode d'ensemble qui combine plusieurs mod\`eles de base
(\textit{base learners}) via un m\'eta-mod\`ele (\textit{meta-learner})
appris sur les pr\'edictions des mod\`eles de base. Contrairement au
bagging ou au boosting, le stacking n'est pas un algorithme unique mais
un cadre g\'en\'eral permettant de combiner des mod\`eles
h\'et\'erog\`enes.
\end{abstract}

\vfill
\tableofcontents
\newpage

% ===================================================================
\section{Principe du stacking}
% ===================================================================

\subsection{Architecture g\'en\'erale}

Le stacking proc\`ede en deux niveaux :

\begin{enumerate}
  \item \textbf{Niveau 0 (base learners)}~: $K$ mod\`eles
    h\'et\'erog\`enes $f_1,\dots,f_K$ sont entra\^\i n\'es
    ind\'ependamment sur les donn\'ees d'apprentissage.
  \item \textbf{Niveau 1 (meta-learner)}~: un m\'eta-mod\`ele $g$ est
    entra\^\i n\'e sur les pr\'edictions des $K$ mod\`eles de base.
\end{enumerate}

Pour une observation $x$, la pr\'ediction finale est~:
\[
\hat y = g\bigl(f_1(x), f_2(x), \dots, f_K(x)\bigr).
\]

\begin{important}
La force du stacking r\'eside dans la diversit\'e des mod\`eles de base.
Chaque apprenant capture des aspects diff\'erents des donn\'ees ; le
m\'eta-mod\`ele apprend \`a corriger leurs erreurs respectives.
\end{important}

\subsection{Pourquoi \c ca marche~?}

\begin{itemize}
  \item \textbf{Compl\'ementarit\'e}~: des mod\`eles de familles
    diff\'erentes (arbres, r\'eseaux, r\'egression lin\'eaire, SVM)
    g\'en\`erent des pr\'edictions diversifi\'ees.
  \item \textbf{Correction des biais}~: le m\'eta-mod\`ele apprend les
    motifs d'erreur des mod\`eles de base.
  \item \textbf{R\'eduction de variance}~: la combinaison de multiples
    estimateurs lisse les pr\'edictions individuelles.
\end{itemize}

\subsection{Proc\'edure de g\'en\'eration des donn\'ees niveau 1}

Pour \'eviter le surapprentissage, les pr\'edictions niveau 1 sont
g\'en\'er\'ees par validation crois\'ee (souvent 5-fold)~:

\begin{enumerate}
  \item Diviser les donn\'ees d'apprentissage $D$ en $V$ plis.
  \item Pour chaque pli $v$~: entra\^\i ner chaque $f_k$ sur
    $D\setminus D_v$, pr\'edire sur $D_v$.
  \item Les pr\'edictions hors-pli pour tous les $K$ mod\`eles forment
    la matrice $Z$ ($N\times K$) d'entr\'ee du m\'eta-mod\`ele.
  \item Entra\^\i ner chaque $f_k$ sur $D$ complet (optionnel).
  \item Entra\^\i ner $g$ sur $(Z, y)$.
\end{enumerate}

% ===================================================================
\section{Impl\'ementation Python avec scikit-learn}
% ===================================================================

\subsection{StackingClassifier}

Scikit-learn propose \texttt{StackingClassifier} et
\texttt{StackingRegressor} depuis la version 0.22.

\begin{lstlisting}
from sklearn.ensemble import StackingClassifier
from sklearn.linear_model import LogisticRegression
from sklearn.tree import DecisionTreeClassifier
from sklearn.svm import SVC
from sklearn.neighbors import KNeighborsClassifier
from sklearn.datasets import load_breast_cancer
from sklearn.model_selection import cross_val_score
import numpy as np

# Donnees
X, y = load_breast_cancer(return_X_y=True)

# Modeles de base heterogenes
base_learners = [
    ('lr',  LogisticRegression(max_iter=5000)),
    ('knn', KNeighborsClassifier(n_neighbors=5)),
    ('svm', SVC(probability=True, random_state=0))
]

# Meta-modele (regression logistique)
meta_learner = LogisticRegression(max_iter=5000)

# Stacking
stack = StackingClassifier(
    estimators=base_learners,
    final_estimator=meta_learner,
    cv=5  # validation croisee 5-fold pour Z
)

# Evaluation
scores = cross_val_score(stack, X, y, cv=5)
print(f"Stacking: {scores.mean():.3f} +/- {scores.std():.3f}")
\end{lstlisting}

\subsection{Param\`etres importants}

\begin{itemize}
  \item \texttt{estimators}~: liste de tuples \texttt{(nom, mod\`ele)}.
  \item \texttt{final\_estimator}~: le m\'eta-mod\`ele (d\'efaut~:
    r\'egression logistique pour classification, Ridge pour
    r\'egression).
  \item \texttt{cv}~: nombre de plis pour g\'en\'erer les
    pr\'edictions niveau 1 (d\'efaut~: 5).
  \item \texttt{stack\_method}~: \texttt{'auto'} (probabilit\'es si
    disponibles, sinon pr\'edictions brutes), \texttt{'predict\_proba'},
    \texttt{'predict'}.
  \item \texttt{passthrough}~: si \texttt{True}, ajoute les
    caract\'eristiques originales $X$ aux pr\'edictions niveau~1.
\end{itemize}

\subsection{StackingRegressor}

\begin{lstlisting}
from sklearn.ensemble import StackingRegressor
from sklearn.linear_model import Ridge, Lasso
from sklearn.tree import DecisionTreeRegressor
from sklearn.ensemble import RandomForestRegressor
from sklearn.datasets import load_diabetes

X, y = load_diabetes(return_X_y=True)

base_learners = [
    ('ridge', Ridge(alpha=1.0)),
    ('lasso', Lasso(alpha=0.1)),
    ('rf',    RandomForestRegressor(n_estimators=100,
                                    random_state=0))
]

stack = StackingRegressor(
    estimators=base_learners,
    final_estimator=Ridge(alpha=1.0),
    cv=5
)

scores = cross_val_score(stack, X, y, cv=5,
                         scoring='neg_mean_squared_error')
print(f"RMSE: {np.sqrt(-scores.mean()):.3f}")
\end{lstlisting}

% ===================================================================
\section{Comparaison avec d'autres méthodes d'ensemble}
% ===================================================================

\begin{tabular}{lccc}
\toprule
M\'ethode & Base learners & Entra\^\i nement & R\'eduction \\
\midrule
Bagging    & m\^eme type & parall\`ele     & variance \\
Boosting   & m\^eme type & s\'equentiel    & biais + variance \\
Voting     & h\'et\'erog\`enes & ind\'ependant & variance \\
Stacking   & h\'et\'erog\`enes & ind\'ependant & biais + variance \\
\bottomrule
\end{tabular}

\subsection{Voting vs Stacking}

Le \texttt{VotingClassifier} combine les pr\'edictions par vote
majoritaire (classification) ou moyenne (r\'egression). Le stacking
apprend \emph{comment} combiner les pr\'edictions, ce qui offre plus de
flexibilit\'e mais risque le surapprentissage.

% ===================================================================
\section{Stratégies avancées}
% ===================================================================

\subsection{Stacking multi-niveaux}

On peut empiler plusieurs couches de m\'eta-mod\`eles~:
\[
\text{Niveau 0} \to \text{Niveau 1} \to \text{Niveau 2}
\]
Attention au surapprentissage~: \`a chaque niveau, les pr\'edictions
doivent \^etre g\'en\'er\'ees par CV.

\subsection{Stacking avec caractéristiques originales}

Avec \texttt{passthrough=True}, le m\'eta-mod\`ele re\c coit les
pr\'edictions \textbf{et} les variables originales $X$, ce qui lui
permet d'apprendre des corrections contextuelles.

\subsection{Feature-weighted stacking}

On peut pond\'erer chaque mod\`ele de base par sa performance estim\'ee
sur un pli de validation plut\^ot que de laisser le m\'eta-mod\`ele
apprendre tous les poids.

% ===================================================================
\section{Exercices}
% ===================================================================

\subsection{Exercice 1 : Stacking manuel}

Impl\'ementer \`a la main un stacking avec 3 classifieurs (KNN,
DecisionTree, SVM) et une r\'egression logistique comme
m\'eta-mod\`ele, sans utiliser \texttt{StackingClassifier}~:

\begin{lstlisting}
from sklearn.model_selection import KFold
import numpy as np

def stacking_manual(X, y, base_learners, meta_learner, n_folds=5):
    kf = KFold(n_splits=n_folds, shuffle=True, random_state=0)
    Z = np.zeros((len(y), len(base_learners)))
    
    for train_idx, val_idx in kf.split(X):
        X_train, X_val = X[train_idx], X[val_idx]
        y_train, y_val = y[train_idx], y[val_idx]
        
        for k, (nom, modele) in enumerate(base_learners):
            modele.fit(X_train, y_train)
            Z[val_idx, k] = modele.predict_proba(X_val)[:, 1]
    
    # Entrainer le meta-modele sur Z
    meta_learner.fit(Z, y)
    
    # Re-entrainer les base learners sur tout X
    for _, modele in base_learners:
        modele.fit(X, y)
    
    return meta_learner, base_learners
\end{lstlisting}

\subsection{Exercice 2 : Comparaison exp\'erimentale}

Comparer les performances de~: (a) chaque mod\`ele de base seul,
(b) VotingClassifier, (c) StackingClassifier sur un jeu de donn\'ees
de classification (breast cancer, iris, ou un jeu synth\'etique).
Interpr\'eter les r\'esultats.

\subsection{Exercice 3 : Stacking r\'egression}

Utiliser \texttt{StackingRegressor} pour combiner Ridge, Lasso et
RandomForest sur les donn\'ees \texttt{diabetes}. Comparer le RMSE
avec chaque mod\`ele individuel.

% ===================================================================
\section{QCM}
% ===================================================================

\noindent
\textbf{1.} Quelle m\'ethode g\'en\`ere les donn\'ees d'entr\'ee du
m\'eta-mod\`ele~?\\
A. Un \'echantillonnage bootstrap\\
B. Une validation crois\'ee\\ 
C. Un tirage al\'eatoire simple\\
D. Une ACP pr\'ealable

\medskip
\noindent
\textbf{2.} Quel est l'avantage principal du stacking par rapport au
VotingClassifier~?\\
A. Il est plus rapide \`a l'entra\^\i nement\\
B. Il apprend \`a combiner les pr\'edictions\\
C. Il utilise moins de m\'emoire\\
D. Il ne n\'ecessite pas de base learners

\medskip
\noindent
\textbf{3.} Quel param\`etre permet d'ajouter les variables
originales au m\'eta-mod\`ele~?\\
A. \texttt{add\_features}\\
B. \texttt{passthrough}\\
C. \texttt{original\_X}\\
D. \texttt{include\_X}

\medskip
\noindent
\textbf{4.} Le stacking r\'eduit principalement~:\\
A. Le biais uniquement\\
B. La variance uniquement\\
C. Le biais et la variance\\
D. Ni le biais ni la variance

\medskip
\noindent
\textbf{5.} Quelle est la principale difficult\'e du stacking
multi-niveaux~?\\
A. Le temps d'ex\'ecution\\
B. Le surapprentissage\\
C. La parall\'elisation\\
D. Le choix des hyper-param\`etres

\medskip
\noindent
\textbf{6.} Que fait \texttt{stack\_method='predict\_proba'}~?\\
A. Utilise les probabilit\'es des mod\`eles de base\\
B. Utilise les classes pr\`edites\\
C. Utilise les \'ecarts-types\\
D. Utilise les r\'esidus

\medskip
\noindent
\textbf{7.} Dans le stacking manuel, pourquoi r\'e-entra\^\i ner les
base learners sur tout $X$ apr\`es la CV~?\\
A. Pour \'economiser de la m\'emoire\\
B. Pour am\'eliorer leurs performances\\
C. Pour r\'eduire le biais\\
D. Ce n'est pas n\'ecessaire

\medskip
\noindent
\textbf{8.} Quel estimateur final est utilis\'e par d\'efaut dans
\texttt{StackingClassifier}~?\\
A. R\'egression logistique\\
B. DecisionTreeClassifier\\
C. SVC\\
D. RandomForestClassifier

\medskip
\noindent
\textbf{9.} Pourquoi les pr\'edictions niveau~1 doivent-elles \^etre
g\'en\'er\'ees par CV~?\\
A. Pour acc\'el\'erer l'entra\^\i nement\\
B. Pour \'eviter le surapprentissage du m\'eta-mod\`ele\\
C. Pour garantir la convergence\\
D. Pour \'equilibrer les classes

\medskip
\noindent
\textbf{10.} Quel type de mod\`eles de base le stacking
peut-il utiliser~?\\
A. Uniquement des arbres\\
B. Uniquement des mod\`eles lin\'eaires\\
C. Des mod\`eles h\'et\'erog\`enes\\
D. Uniquement des SVM

\newpage
\section*{Corrig\'e du QCM}

\medskip
1-B, 2-B, 3-B, 4-C, 5-B, 6-A, 7-B, 8-A, 9-B, 10-C.

\end{document}
